\clearpage
\thispagestyle{plain}
\begin{center}
{\fontsize{16}{20}\bfseries Abstract\par}
\end{center}
\vspace{0.1in}

Organisations and individuals increasingly hold their knowledge in documents: reports,
spreadsheets, scanned files, slide decks and source code. Finding a precise answer inside
such a collection, or turning it into a chart, a summary or a clear picture of how its ideas
connect, still demands a great deal of manual reading. General-purpose chatbots help, but
they cannot ground their answers in a private library, they invent facts when evidence is
thin, and they cannot read a scanned page, draw a chart from a table, or show how the
concepts in a corpus relate to one another.

This thesis presents Synapse, a multi-tenant platform that lets a user connect a library
of documents and then understand, question, visualise and map that knowledge through a
single web application. Synapse ingests each document through a pipeline that detects
layout, reads born-digital and scanned pages (the latter through a vision-language model),
chunks and embeds the text, and clusters it for topical coverage. A retrieval-augmented
chat answers questions using a corrective, multi-agent retrieval loop that grades its own
evidence and abstains when the library cannot support an answer, citing the exact source
passage for every claim. An agent mode turns the same libraries into artefacts: it plans
its work, asks the user a clarifying question when a request is genuinely ambiguous, and
produces interactive charts, written documents, PDFs and generated images. A separate
knowledge-graph module extracts the entities and relationships from a processed library
and draws them as an interactive graph that the user can explore. Teams share libraries
and chat together in real time.

The platform was evaluated on the DOUBLE-BENCH document benchmark for retrieval and
answer quality, on a set of purpose-built spreadsheet, Python and C++ datasets for
structured-document understanding, and through functional testing of the agent and graph
features. The results show accurate document-level retrieval, well-grounded and honest
answers, and reliable artefact generation, demonstrating that a disciplined pipeline paired
with a corrective, multi-agent retrieval design can deliver practical document intelligence
on a user's own private collection.

\vspace{12pt}
\noindent\textbf{Keywords:} Retrieval-Augmented Generation, Document Intelligence,
Multi-Agent Systems, Corrective RAG, Knowledge Graphs, Vision-Language Models,
Question Answering, Data Visualisation, Large Language Models.
